\documentclass[10pt]{article}
\usepackage[utf8]{inputenc}
\usepackage[T1]{fontenc}
\usepackage{amsmath}
\usepackage{amsfonts}
\usepackage{amssymb}
\usepackage[version=4]{mhchem}
\usepackage{stmaryrd}

\title{A SHORT PROOF OF CAUCHY'S POLYGONAL NUMBER THEOREM }

\author{MELVYN B. NATHANSON}
\date{}


%New command to display footnote whose markers will always be hidden
\let\svthefootnote\thefootnote
\newcommand\blfootnotetext[1]{%
  \let\thefootnote\relax\footnote{#1}%
  \addtocounter{footnote}{-1}%
  \let\thefootnote\svthefootnote%
}

%Overriding the \footnotetext command to hide the marker if its value is `0`
\let\svfootnotetext\footnotetext
\renewcommand\footnotetext[2][?]{%
  \if\relax#1\relax%
    \ifnum\value{footnote}=0\blfootnotetext{#2}\else\svfootnotetext{#2}\fi%
  \else%
    \if?#1\ifnum\value{footnote}=0\blfootnotetext{#2}\else\svfootnotetext{#2}\fi%
    \else\svfootnotetext[#1]{#2}\fi%
  \fi
}

\begin{document}
\maketitle


\begin{abstract}
This paper presents a simple proof that every nonnegative integer is the sum of $m+2$ polygonal numbers of order $m+2$.
\end{abstract}

Let $m \geqslant 1$. The polygonal numbers of order $m+2$ are the integers

$$
p_{m}(k)=\frac{m}{2}\left(k^{2}-k\right)+k
$$

for $k=0,1,2, \ldots$. Fermat [3] asserted that every nonnegative integer is the sum of $m+2$ polygonal numbers of order $m+2$. For $m=2$, Lagrange [5] proved that every nonnegative integer is the sum of four squares $p_{2}(k)=k^{2}$. For $m=1$, Gauss [4] proved that every nonnegative integer is the sum of three triangular numbers $p_{1}(k)=\left(k^{2}+k\right) / 2$, or, equivalently, that every positive integer $n \equiv 3(\bmod 8)$ is the sum of three odd squares. Cauchy [1] proved Fermat's statement for all $m \geqslant 3$, and Legendre [6] refined and extended this result. For $m \geqslant 3$ and $n \leqslant 120 m$, Pepin [8] published tables of explicit representations of $n$ as a sum of $m+2$ polygonal numbers of order $m+2$, at most four of which are different from 0 or 1. Dickson [2] prepared similar tables. Pall [7] obtained important related results on sums of values of a quadratic polynomial.

Uspensky and Heaslet [9, p. 380] and Weil [10, p. 102] have written that there is no short and easy proof of Cauchy's polygonal number theorem. The object of this note is to present a short and easy proof.

Because of Pepin's and Dickson's tables, it suffices to consider only $n \geqslant 120 m$. For completeness, I also include a proof of Cauchy's lemma.

Cauchy's Lemma. Let $a$ and $b$ be odd positive integers such that $b^{2}<4 a$ and $3 a<b^{2}+2 b+4$. Then there exist nonnegative integers $s, t, u, v$ such that


\begin{align*}
& a=s^{2}+t^{2}+u^{2}+v^{2}  \tag{1}\\
& b=s+t+u+v \tag{2}
\end{align*}


Proof. Since $a$ and $b$ are odd, it follows that $4 a-b^{2} \equiv 3(\bmod 8)$, and so, by Gauss's triangular number theorem, there exist odd integers $x \geqslant y \geqslant z>0$ such that


\begin{equation*}
4 a-b^{2}=x^{2}+y^{2}+z^{2} . \tag{3}
\end{equation*}


\footnotetext{Received by the editors August 26, 1985 and, in revised form, December 20, 1985.\\
1980 Mathematics Subject Classification (1985 Revision). Primary 11P05, 11B83, 11B13.\\
Key words and phrases. Polygonal numbers, sums of squares, Cauchy-Fermat theorem.
}Choose the sign of $\pm z$ so that $b+x+y \pm z \equiv 0(\bmod 4)$. Define integers $s, t, u, v$ by

$$
\begin{gathered}
s=\frac{b+x+y \pm z}{4}, \quad t=\frac{b+x}{2}-s=\frac{b+x-y \mp z}{4}, \\
u=\frac{b+y}{2}-s=\frac{b-x+y \mp z}{4}, \quad v=\frac{b \pm z}{2}-s=\frac{b-x-y \pm z}{4} .
\end{gathered}
$$

Then equations (1) and (2) are satisfied, and $s \geqslant t \geqslant u \geqslant v$. To show these integers are nonnegative, it suffices to prove that $v \geqslant 0$, or $v>-1$. This is true if $b-x-y$ $-z>-4$, or, equivalently, if $x+y+z<b+4$. The maximum value of $x+y+z$ subject to the constraint (3) is $\sqrt{12 a-3 b^{2}}$, and the inequality $3 a<b^{2}+2 b+4$ implies that $x+y+z \leqslant \sqrt{12 a-3 b^{2}}<b+4$. This proves the lemma.

Theorem 1. Let $m \geqslant 3$ and $n \geqslant 120 m$. Then $n$ is the sum of $m+1$ polygonal numbers of order $m+2$, at most four of which are different from 0 or 1.

Proof. Let $b_{1}$ and $b_{2}$ be consecutive odd integers. The set of numbers of the form $b+r$, where $b \in\left\{b_{1}, b_{2}\right\}$ and $r \in\{0,1, \ldots, m-3\}$, contains a complete set of residue classes modulo $m$, and so $n \equiv b+r(\bmod m)$ for some $b \in\left\{b_{1}, b_{2}\right\}$ and $r \in\{0,1, \ldots, m-3\}$. Define


\begin{equation*}
a=2\left(\frac{n-b-r}{m}\right)+b=\left(1-\frac{2}{m}\right) b+2\left(\frac{n-r}{m}\right) . \tag{4}
\end{equation*}


Then $a$ is an odd integer, and


\begin{equation*}
n=\frac{m}{2}(a-b)+b+r . \tag{5}
\end{equation*}


If $0<b<\frac{2}{3}+\sqrt{8(n / m)-8}$, then the quadratic formula implies that

$$
b^{2}-4 a=b^{2}-4\left(1-\frac{2}{m}\right) b-8\left(\frac{n-r}{m}\right)<0
$$

and so $b^{2}<4 a$. Similarly, if $b>\frac{1}{2}+\sqrt{6(n / m)-3}$, then $3 a<b^{2}+2 b+4$. Since the length of the interval


\begin{equation*}
I=\left(\frac{1}{2}+\sqrt{6\left(\frac{n}{m}\right)-3}, \frac{2}{3}+\sqrt{8\left(\frac{n}{m}\right)-8}\right) \tag{6}
\end{equation*}


is greater than 4, it follows that $I$ contains two consecutive odd positive integers $b_{1}$ and $b_{2}$. Thus, there exist odd positive integers $a$ and $b$ that satisfy (5) and the inequalities $b^{2}<4 a$ and $3 a<b^{2}+2 b+4$. Cauchy's Lemma implies that there exist $s, t, u, v$ satisfying (1) and (2), and so

$$
\begin{aligned}
n & =\frac{m}{2}(a-b)+b+r=\frac{m}{2}\left(s^{2}-s\right)+s+\cdots+\frac{m}{2}\left(v^{2}-v\right)+v+r \\
& =p_{m}(s)+p_{m}(t)+p_{m}(u)+p_{m}(v)+r .
\end{aligned}
$$

This completes the proof.\\[0pt]
Note that this result is slightly stronger than Cauchy's theorem. Legendre [6] proved that every sufficiently large integer is the sum of five polygonal numbers of order $m+2$, one of which is either 0 or 1. This can also be easily proved.

Theorem 2. Let $m \geqslant 3$. If $m$ is odd, then every sufficiently large integer is the sum of four polygonal numbers of order $m+2$. If $m$ is even, then every sufficiently large integer is the sum of five polygonal numbers of order $m+2$, one of which is either 0 or 1.

Proof. There is an absolute constant $c$ such that if $n>c m^{3}$, then the length of the interval $I$ defined in (6) is greater than $2 m$, and so $I$ contains at least $m$ consecutive odd integers.

If $m$ is odd, these form a complete set of residues modulo $m$, and so $n \equiv b$ $(\bmod m)$ for some odd number $b \in I$. Let $r=0$. Define $a$ by formula (4).

If $m$ is even and $n>c m^{3}$, then $n \equiv b+r(\bmod m)$ for some odd integer $b \in I$ and $r \in\{0,1\}$. Define $a$ by (4).

In both cases, the theorem follows immediately from Cauchy's Lemma.

\section*{References}
\begin{enumerate}
  \item A. Cauchy, Démonstration du théorème général de Fermat sur les nombres polygones, Mém. Sci. Math. Phys. Inst. France (1) 14 (1813-15), $177-220=$ Oeuvres (2), vol. 6, 320-353.
  \item L. E. Dickson, All positive integers are sums of values of a quadratic function of $x$, Bull. Amer. Math. Soc. 33 (1927), 713-720.
  \item P. Fermat, quoted in T. L. Heath, Diophantus of Alexandria, Dover, New York, 1964, p. 188.
  \item C. F. Gauss, Disquisitiones arithmeticae, Yale Univ. Press, New Haven, Conn., and London, 1966.
  \item J. L. Lagrange, Démonstration d'un théorème d'arithmétique, Nouveaux Mémoires de l'Acad. Royale des Sci. et Belles-L. de Berlin, 1770, pp. 123-133 = Oeuvres, vol. 3, pp. 189-201.
  \item A.-M. Legendre, Théorie des nombres, 3rd ed., vol. 2, 1830, pp. 331-356.
  \item G. Pall, Large positive integers are sums of four or five values of a quadratic function, Amer. J. Math. 54 (1932), 66-78.
  \item T. Pepin, Démonstration du théorème de Fermat sur les nombres polygones, Atti Accad. Pont. Nuovi Lincei 46 (1892-93), 119-131.
  \item J. V Uspensky and M. A. Heaslet, Elementary number theory, McGraw-Hill, New York and London, 1939.
  \item A. Weil, Number theory, an approach through history from Hammurabi to Legendre, Birkhäuser, Boston, Mass., 1983.
\end{enumerate}

Department of Mathematics, Rutgers University, Newark, New Jersey 07102\\
Current address: Office of the Provost and Vice President for Academic Affairs, Lehman College (CUNY), Bronx, New York 10468


\end{document}